\section[\appendixname~\thesection]{Validation Methods}
\label{sec:appendix_c}

Suppose the training set contains \(n\) samples, \(x_1, x_2, \dots, x_n\).
Let
\[
	\mathcal{M} : \mathcal{D}_{\mathrm{train}} \times \mathcal{D}_{\mathrm{val}} \to \mathbb{R}
\]
denote an evaluation function that maps a training set and a validation set to a scalar performance metric.

The following defines the time-series cross-validation procedure used in this study.

Let \(K\) denote the number of validation folds.
We require the samples to be ordered chronologically by their observation times.
That is, for any \(1 \le i < j \le n\), sample \(x_i\) is observed before sample \(x_j\).
The ordered samples are partitioned into \(K + 1\) consecutive blocks, denoted by \(B_1, B_2, \dots, B_K, B_{K + 1}\), with approximately equal sizes.
Each block must contain at least one sample; the first block serves as the initial training window.

For each validation fold \(B_i\), where \(i = 2, \dots, K+1\), the model is trained using all preceding blocks, \(\bigcup_{j=1}^{i-1} B_j\), and evaluated on \(B_i\).
The mean performance metric across the \(K\) validation folds is defined as
\[
	\bar{m} = \frac{1}{K} \sum_{i=2}^{K + 1} \mathcal{M} \left( \bigcup_{j=1}^{i-1} B_j, B_i \right).
\]

The implementation uses \texttt{time\_series\_cross\_validation}, which partitions observation indices with \texttt{numpy.array\_split}; we use \(K=4\) validation folds.

The following defines the repeated random holdout validation procedure used in this study.

Let \(T\) denote the number of iterations and \(r\) denote the validation rate.
For each iteration \(t = 1, 2, \dots, T\), let
\[
	u = \max \left(2, \left\lfloor nr \right\rfloor \right)
\]
denote the number of samples selected for validation.

Randomly select \(u\) distinct indices \(i_1^{(t)}, i_2^{(t)}, \dots, i_u^{(t)}\).
Let
\[
	I_t = \left\{ i_1^{(t)}, i_2^{(t)}, \dots, i_u^{(t)} \right\}
\]
denote the set of validation indices for iteration \(t\).
The performance metric for iteration \(t\) is defined as
\[
	m_t = \mathcal{M} \left( \{x_j : j \notin I_t\}, \{x_j : j \in I_t\} \right).
\]

The mean performance metric across all \(T\) iterations is then defined as
\[
	\bar{m} = \frac{1}{T} \sum_{t=1}^{T} m_t.
\]

% In the implementation, the default validation rate is \(r=0.2\) and \(T\) defaults to 4.
The interpolation test set contains only interior observations, with the first and last observations retained in the training set.
During hyperparameter selection, however, validation masks are sampled from all observations in the training set, including its endpoints.
This retains more possible validation subsets in very small training sets, but some validation splits may involve extrapolation.
The tuning results therefore reflect a mixture of interpolation and extrapolation performance, whereas the final test evaluates only interpolation.
